\chapter[In-Game User Interface Implementation \textsuperscript{Novotny}]{In-Game User Interface Implementation}
\label{chap:In-Game User Interface}

\section{Organization of this Chapter}
This chapter treats the implementation of the in-game user interface. A section has been devoted to each of three display options the application provides:

\begin {itemize}
	\item {Piano Roll}
	\item {Score View}
	\item {Tablature View}
\end {itemize}

The first section gives an overview of the architechture behind the in-game user interface.

\section{Design}
The implementation of all three display options is based on the {\tt QQuickItem} drawing approach explained in the \textit{Technologies} chapter (\ref{chap:Technologies}). 

The three display options require many common resources. Application logic and calculations used in for the individual display options are very similar in a number of ways. The class hierachy used for implementing the display options helps avoid redundancies while maintaining flexibility.

\subsection{Visuals}
{\tt Visuals} handles tasks relevant to all display options. These tasks are described throughout the following subsections.

\subsubsection{Providing Data}
{\tt Visuals} provides a property {\tt game} of type {\tt MainGame}.
This property is used for accessing:
\begin {itemize}
	\item {tracks of the score}
	\item {track played by the user}
\end {itemize}

\subsubsection{Scrolling}
\subsubsection{Determining the appropriate Time Span to display}
In order to provide relevant instructions and feedback to the user, the amount of information presented needs to be adjusted to the complexity of the score. Primarily this means displaying the right amount of notes simultaneously.
{\tt Visuals} provides the attribute {\tt msPerScreen} which sets the time span to be displayed on the screen. This attribute is set to an appropriate value depending on the density of notes in the track. A value of 1000 means that one second of music will be displayed. 


\subsection{QQI\_Based\_Visuals}
The class {\tt QQI\_Based\_Visuals} extends {\tt Visuals}. This class offers a range of functions relevant for drawing using the {\tt QtQuickItem} based approach.
The reason for introducing this second subclass is maintaining flexibility in terms of how the components are painted. While experience has shown that using the canvas-based drawing approach is not suitable for drawing complex real-time oriented components due to limitations of the current rendering implementation this architecture allows quick adoptions in the future.
Functions provided by {\tt QQI\_Based\_Visuals} are discussed throughout the following subsections.

\subsubsection{Converting Time Stamps to Positions}
Objects such as:
\begin{itemize}
	\item{notes}
	\item{bar lines}
	\item{bending profiles}
\end{itemize}
are a associated with a timestamp relative to the start of the song.
The function {\tt calculatePosition} has a number of overloads. It accepts either a time stamp or a {\tt Note} object and a boolean flag indicating whether or not the object is contained in the score (as opposed to played by the user). {\tt calculatePosition} returns the horizontal position of objects taking the following factors into account:
\begin{itemize}
	\item{scaling of the component}
	\item{horizontal shifts that occur if the song is started at a position other than its beginning (refer to the chapter on \textit{Design} for further details)}
\end{itemize}

\subsubsection{Loading Textures}
Textures such as:
\begin{itemize}
	\item{triangle textures (pointing left and right) that indicate timing mistakes}
	\item{digit textures required for printing the score above notes}
\end{itemize}
are required by all display options. The function {\tt loadTextures} loads these textures.
Note that every display option implements their on {\tt loadTextures} function which calls the {\tt QQI\_Based\_Visuals::loadTextures} and loads add loads additionals textures it requires.

\subsubsection{Drawing Bar Lines}
Drawing bar lines significantly increases the readability of scores and is used every display option.
{\tt QQI\_Based\_Visuals} offers the function {\tt getBarsNode}. The function returns a {\tt QSGNode} containing at least all bar nodes in the visible scope. This node can be added to the parent node by the calling display option.

The start time of every bar is contained in the {\tt Score} object maintained by the {\tt game} object. {\tt getBarsNode} uses the function {\tt calculatePosition} to position the bar lines accordingly.

\subsubsection{Displaying Target Notes and Played Notes}
The function {\tt getTargetNotes} returns a node containing all target notes of the selected track. 
Similarly the function {\tt getPlayedNotes} returns a node containing all notes played by the user.

In order to improve performance only notes within the visible time range are added to the node.
The variable {\tt startIndex} stores the index of the first note within the track that needs to be drawn. This variable is updated continuously.

The functions use the virtual function {\tt getNodeForNote} that is implemented by every display option. This allows the display option to define the appearance of the notes while avoiding redundancies.

\subsubsection{Score for Target Notes}
As explained in \textit{User Interaction Concept} (\ref{chap:UserInteractionConcept}) every note of the track the user is trying to play is associated with a numeric score. The maximum score for every note is 50 points, deductions are made for playing the wrong pitch and timing mistakes.

By default every target note has a score of -1. This value is updated by the {\tt Evaluator} as soon as a played note is associated with the target note. If one of the display options recognizes a target note with a non-negative score the numeric value is printed above the target note.

The function {\tt getNumberNode(int number, int size, int left, int top)} returns a {\tt QSGNode} that displays the specified number and has the specified position and size.

\subsubsection{Drawing Text Nodes}
Qt does not offer a mechanism for drawing text nodes. Drawing numeric values representing the score of a note requires a number of steps. 

The funtion {\tt loadTextures} creates a {\tt QList} of textures containing a texture for every digit between 0 and 9.
This is done by loading an image containing all digits and splitting the image into ten equally sized pieces using the function {\tt createSubImage}.

\begin{lstlisting}
	QImage QQI_Based_Visuals::createSubImage(QImage* image, const QRect & rect)
	{
	    size_t offset = rect.x() * image->depth() / 8
	                    + rect.y() * image->bytesPerLine();
	    return QImage(image->bits() + offset, rect.width(), rect.height(),
	                  image->bytesPerLine(), image->format());
	}
\end{lstlisting}

The function {\tt getNumberNode} creates the node using the following steps:

\begin{enumerate}
	\item{split the number into a list of digits}
	\item{create a parent node that will contain all digit nodes}
	\item{create a {\tt QSGSimpleTextureNode} for every digit using the textures created in {\tt loadTextures} and add it to the parent node}
\end{enumerate}

\subsection{Drawing notes in the score view}
Drawing notes in the score view is special because the {\tt QSGTexture}s are 
generated from Scalable Vector Graphic (SVG) files. SVG files are a form of picture format where, instead of saving an array of pixels in the file, the drawing is described in XML in the form of shapes, colors or gradients. To generate a texture from an SVG file several steps have to be made. The following sample code shows a simple way to generate textures.
\begin{lstlisting}
 QSvgRenderer svgRenderer(texturePath);
 QImage image(width, height, QImage::Format_ARGB32);
 image.fill(0xFFFFFF);
 QPainter painter(&image);
 painter.setRenderHint(QPainter::Antialiasing);
 svgRenderer.render(&painter);
 QSGTexture *texture = window()->createTextureFromImage(image);
\end{lstlisting}

The {\tt QImage} object is the temporary object that will store a rendered version of the SVG file. {QSvgRenderer} uses the provided {\tt QPainter} to paint the SVG document into the {\tt QImage}. After that {\tt createTextureFromImage()} can be used to generate a texture from the QImage for use in the scene graph.
\\
This type of texture is used for the following types of objects:
\begin{enumerate}
\item note textures in black, green and red
\item dotted notes
\item sharp (\#) sign
\item grace notes
\item clef
\end{enumerate}

\subsection{Painting Process for all Display Options}
Ever display options implements the function {\tt updatePaintNode} defined in {\tt QQuickItem}. The drawing process implemented in these functions generally follows the following flow:
\begin{enumerate}
	\item{
		Preparations such as:
		\begin{itemize}
			\item{deleting all contents of the scene graph and creating a new parent node (refer to the chapter \textit{Technologies} (\ref{chap:Technologies}))}
 			\item{retrieving the current game time}
			\item{updating the scroll bar position}
		\end{itemize}
		Where possible these preparations are done by functions of {\tt QQI\_Based\_Visuals} in a generic way in order to avoid code duplicates.
	}
	\item{
		Drawing objects specific to the display option such as:
		\begin{itemize}
			\item{a piano graphic and coloured lines in {\tt PianoRoll}}
			\item{strings in {\tt TablatureView}}
			\item{clef in {\tt ScoreView}}
		\end{itemize}
	}
	\item{
		Drawing objects required by all display options such as:
		\begin{itemize}
			\item{a vertical line indicating the current point in time}
			\item{bar lines}
		\end{itemize}
		The correpsonding nodes for these objects are created by {\tt QQI\_Based\_Visuals}.
	}
	\item{
		Drawing target notes and played notes. This process has been covered in the previous section.
	}
\end{enumerate}
